
\subsection{A class-group gate for the septic residual source}

\begin{lemma}[Split-at-three ray classes]
\label{lem:lm-septic-ray-gate}
Let $K/\mathbb Q$ be Galois and contain $\mu_3$.
Suppose it has one prime $\mathfrak p$ above $3$, with
$(3)=\mathfrak p^e$ and $3\nmid e$, and exactly two primes
$\mathfrak q_1,\mathfrak q_2$ above $7$, each of residue
degree two.  Put $S=\{v:v\mid21\}$.
Then
\[
 \ker\bigl(H^1(G_{K,S},\mathbb F_3)
       \longrightarrow H^1(K_{\mathfrak p},\mathbb F_3)\bigr)
 \simeq\operatorname{Hom}(\operatorname{Cl}(K),\mathbb F_3).
\]
In particular, the local group at $3$ surjects onto the
maximal pro-$3$ extension unramified outside $S$ if and
only if $\operatorname{Cl}(K)[3]=0$.
\end{lemma}
\begin{proof}
A character in the kernel cuts out a cyclic cubic extension
unramified at $3$, split at $\mathfrak p$, and unramified
outside the two primes over $7$.  Its possible ramification
there is tame, so its conductor divides
$\mathfrak m=\mathfrak q_1\mathfrak q_2$.
Class field theory identifies these characters with the
characters of the ray class group modulo $\mathfrak m$,
quotiented by the class of $\mathfrak p$.
The relevant exact sequence is
\[
 \mathcal O_K[1/3]^\times\longrightarrow
 (\mathcal O_K/\mathfrak m)^\times\longrightarrow
 C_{\mathfrak m}^{\{\mathfrak p\}}\longrightarrow
 \operatorname{Cl}(\mathcal O_K[1/3])\longrightarrow0.
\]
Each residue field is $\mathbb F_{49}$, whose multiplicative
group has three-primary part of order three.  The two
primes restrict to the distinct primes of
$\mathbb Q(\zeta_3)$ above $7$: otherwise the Galois action
on primes and the existence of both restrictions would
give more than two primes in $K$.
Choose the common generator $2$ of the order-three subgroup
in $\mathbb F_7^\times$ inside each residue field.
The three-primary projection of $\zeta_3$ has coordinate
vectors $(1,2)$ or $(2,1)$; that of the rational
$S$-unit $3$ has equal nonzero coordinates, because
$3^{16}=4$ in $\mathbb F_7$.
These two vectors are linearly independent over $\mathbb F_3$.
Thus the images of $\mu_3$ and $3$ generate the entire
three-primary residue group.
The three-primary part of $C_{\mathfrak m}^{\{\mathfrak p\}}$
therefore equals that of
$\operatorname{Cl}(\mathcal O_K[1/3])$.
Finally $[\mathfrak p]$ has order dividing $e$, which is
prime to three.  Removing $\mathfrak p$ does not change
the three-primary class group, proving the isomorphism.
The last assertion is the Frattini criterion for a
pro-$3$ group, as in the quintic residual-source proof.
\end{proof}

\begin{proposition}[Residual data and the class-group reduction]
\label{prop:lm-septic-residual-gate}
Let $A_7=\operatorname{Jac}(v^2+v=u^7)$,
$F=\mathbb Q(\zeta_7)$, and $K=\mathbb Q(A_7[3])$.
Then $K$ satisfies the hypotheses of
Lemma~\ref{lem:lm-septic-ray-gate}, with $e=56$.
More precisely,
\[
 [K:\mathbb Q]=336,\quad\operatorname{Gal}(K/F)=C_{56},
 \quad |D_K|=3^{330}7^{380}.
\]
At $3$ there is one prime with $(e,f)=(56,6)$.
At $7$ there are two primes, each with $(e,f)=(84,2)$.
If $H=C_7\subset C_{56}$ and $K_0=K^H$, then
\[
 [K_0:\mathbb Q]=48,\qquad |D_{K_0}|=3^{42}7^{44},
 \qquad\operatorname{Cl}(K_0)[3]=0.
\]
The $H$-invariant part of the localization kernel for $K$
is zero.  Its possible remaining part has
$\mathbb F_3$-dimension divisible by six; its vanishing
is exactly the condition
$\operatorname{Cl}(K)[3]=0$.
\end{proposition}
\begin{proof}
The same CM endomorphism argument as for $A_5$ gives
$F,\mathbb Q(\zeta_3)\subset K$.
The prime $3$ is inert in $F$, and the residual module is
one dimensional over $\mathbb F_{729}$.
Polarization bounds the $G_F$-image by
$\{a:a^{28}\in\mathbb F_3^\times\}=C_{56}$.
The differential type is $\{1,2,3\}$, with reflex
$\{1,4,5\}$; local CM reciprocity at the good prime
$3$ gives the power $-(1+3^4+3^5)=-325$.
Since $\gcd(325,728)=13$, its image has order $56$.
Thus local inertia fills the global $G_F$-image, and
the local decomposition group at $3$ has order
$56\cdot6=336$, the whole global group.

\citet[Proposition~3.12(2), author p.~16]{ShuExpanded2021}, applied
to $(r,s,t,\delta)=(1,1,5,1)$, gives conductor $\lambda_7^2$.
His Equation~(3.14) on that page expresses the local
unit character as the product of a seventh-power
Hilbert-symbol character and a quadratic character of
conductor one.  The exact conductor two makes the
order-seven factor nontrivial.  Thus the unit image is
$\mu_{14}$, retaining its order modulo three, and
$e_7(K/F)=14$.

For the residue degree, take $\lambda=1-\zeta_7$.
The global reciprocity formula on the principal idele
$\lambda$ shows that its local character at $7$, modulo
inertia, is the reflex product
\[
 \alpha=(1-\zeta_7)(1-\zeta_7^4)(1-\zeta_7^5)
       \quad\text{in }\mathbb F_{729}^{\times},
\]
up to inversion and coefficient conjugation.
Indeed all other finite unit factors are trivial by
good reduction, while the local-three factor is the
inverse reflex norm.
In the field
$\mathbb F_3[z]/(z^6+z^5+z^4+z^3+z^2+z+1)$, direct
reduction gives
\[
 \alpha^2=-z^3,\qquad \alpha^{14}=-1,\qquad
 \alpha^{28}=1.
\]
Hence its class modulo $\mu_{14}$ has order two.
Therefore $f_7=2$, and the number of primes is
$336/(84\cdot2)=2$.
In particular the prime above $7$ is \emph{not} unique.

The conductor--discriminant calculation has fifty-five
nontrivial characters of conductor exponent one at
$3$, whose prime norm in $F$ is $3^6$.
At $7$, forty-eight characters have conductor exponent
two, four have exponent one, and three nontrivial
characters are unramified.  As $|D_F|=7^5$, this gives
\[
 |D_K|=7^{5\cdot56}\,3^{6\cdot55}\,7^{2\cdot48+4}
       =3^{330}7^{380}.
\]
The corresponding quotient $K_0/F$ has degree eight,
inertia of order two at $7$, and gives
$|D_{K_0}|=3^{42}7^{44}$.
The stronger class-number-one assertion in
Corollary~\ref{cor:septic-auxiliary-class-one} below proves
the required three-class vanishing, using only this
discriminant and the explicit-formula test function used above.

The field $K_0$ has the same two residue-degree-two
primes above $7$ and one prime of ramification index
eight above $3$.  Lemma~\ref{lem:lm-septic-ray-gate}
therefore makes its entire split-at-three localization
kernel zero.  Inflation--restriction through the
degree-seven extension $K/K_0$, together with local
restriction of degree seven, identifies this with the
$H$-invariant kernel for $K$.
Finally $\Phi_7$ is irreducible over $\mathbb F_3$,
since $3$ has order six modulo seven.  All nontrivial
$\mathbb F_3[C_7]$-constituents have dimension six.
The earlier ray-class lemma identifies their possible
sum with the remaining three-class-group gate.
\end{proof}

\begin{corollary}[The auxiliary class number]\label{cor:septic-auxiliary-class-one}
For the degree-$48$ field $K_0$ above, one has $h(K_0)=1$.
\end{corollary}
\begin{proof}
If its class number were greater than one, its Hilbert class
field would be a totally complex field of degree $d\ge96$
with the same root discriminant
$3^{7/8}7^{11/12}$. Use the kernel of
Proposition~\ref{prop:quintic-three-class} with $b=51/5$,
and discard every nonnegative prime contribution. The
unconditional explicit formula gives
\[
 \operatorname{rd}\ge B\exp\!\left(-\frac{136}{5d}\right),
 \qquad d\ge96.
\]
The endpoint and tail bounds for $B$ are exactly those already
specified there. Re-evaluating that integral with $192$-bit
ball arithmetic gives the rational certified inequalities
\[
 B\exp(-17/60)>\frac{1585}{100},
 \qquad 3^{7/8}7^{11/12}<\frac{1557}{100}.
\]
The lower bound increases with $d$. The incompatible bounds
exclude the Hilbert class field and prove the assertion.
The calculation does not compute the class groups of $K$ or $M$.
\end{proof}

\begin{remark}[Why the smaller field is used]
\label{rem:septic-smaller-field}
The root discriminant of $K$ is approximately $26.5685$, larger
than the unconditional asymptotic complex-field discriminant threshold.
The quintic discriminant argument therefore does not exclude a
nontrivial $C_7$-constituent by its dimension alone. We will first
control the whole geometric action, identify the field, and reduce
its three-class group to a degree-$56$ field. Principal prime ideals
in that smaller field then supply the stronger explicit-formula bound
needed to prove vanishing.
\end{remark}
